\documentclass[10pt,a4paper]{amsart}
\usepackage[margin=19mm]{geometry}
\usepackage{amsmath,amssymb,mathtools}
\usepackage[scr=rsfs,cal=euler]{mathalfa}
\usepackage{xfrac}
\usepackage[unicode,hidelinks,bookmarks=false]{hyperref}
\newcommand{\ie}{i.e.\ }
\newcommand{\cf}{cf.\ }
\newcommand{\resp}{respectively\ }
\newcommand{\Subsection}{\subsection}
\providecommand{\nobreakdash}{\nobreak\hbox{-}}
\setlength{\emergencystretch}{2em}
\multlinegap0pt
\usepackage{amssymb}
\usepackage[shortlabels]{enumitem}
\usepackage{tikz}
\usepackage{bm}
\usepackage{mathtools}
\newtheorem{proposition}{Proposition}
\newtheorem{lemma}{Lemma}
\usepackage{cleveref}
\crefname{proposition}{Proposition}{Propositions}
\crefname{lemma}{Lemma}{Lemmas}
\newcommand{\cherry}{\tikz[baseline=-0.15ex]{\draw[darkgray, thick] (0,0)--(0.1,0.17320508075)--(0.2,0); \draw[darkgray, fill=white] (0,0) circle (0.05); \draw[darkgray, fill=white] (0.1,0.17320508075) circle (0.05); \draw[darkgray, fill=white] (0.2,0) circle (0.05);}}
\newcommand{\cl}[1]{\langle #1 \rangle}
\newcommand{\gen}[1]{\langle #1 \rangle_\pm}
\newcommand{\sset}[1]{\left\{{#1}\right\}}
\title{Classification of equiangular lines\\with fixed angle $\arccos(1/(1+2\sqrt2))$}
\author{Theodore Gossett, Zilin Jiang, Adam Teets, Zoe Wellner}
\date{Source: arXiv:2603.02469v1 (CC BY 4.0)}
\begin{document}
\maketitle

\noindent\emph{Source and licence.} Classification of equiangular lines\\with fixed angle $\arccos(1/(1+2\sqrt2))$. Version arXiv:2603.02469v1 (CC BY 4.0), distributed by arXiv under the Creative Commons Attribution 4.0 International licence, http://creativecommons.org/licenses/by/4.0/, as recorded in the arXiv licence metadata for this exact version and shown on the abstract page https://arxiv.org/abs/2603.02469v1. Full licence notice: \texttt{licenses/arxiv-2603.02469-CC-BY-4.0.txt}.\par\smallskip

\noindent\textit{Editorial scope.} Counted unit: the article's lemma lem:min-forb-switching (source0.tex lines 274--276). The article's complete original proof of it is delivered with it (source0.tex lines 278--284, one \texttt{proof} environment of 7 lines). No mathematics is written, repaired or re-derived: the statement and the proof are the article's own lines, in the article's own notation, and no retained line is substituted or re-flowed.  Retained uncounted as the source-local context the proof needs: lines 256--258; lines 264--266.  Nothing was omitted from any retained span, so the package carries no omission record.  No retained line contains a citation, so the wrapper carries no bibliography and no bibliography entry.  No retained line contains a figure environment, an image-inclusion command or drawing code, so no image is delivered and the package declares no asset dependency.  Editorial formatting only, in addition: an AMS \texttt{amsart} preamble that restates the article's own declarations and macros used by the retained lines, namely the environments \texttt{proposition}, \texttt{lemma} on separate counters, together with the standard packages those lines need.  The retained environments print in this document as Proposition 1, Lemma 1 in order of appearance, whereas the article prints the same items with the numbering of its own full document.  The complete native source of the article as distributed in the arXiv e-print for this exact version is bundled unchanged as \texttt{sources/195621/source0.tex}.
\begin{proposition} \label{lem:uniq}
    For every $G \in \cl\cherry$ of order at least $7$, and every $U \subseteq V(G)$, if the switching $G^U$ of $G$ with respect to $U$ is also in $\cl\cherry$, then $U$ is trivial, that is, $U = \varnothing$ or $U = V(G)$.
\end{proposition}

\begin{proposition} \label{lem:min-forb}
    Any minimal forbidden subgraph for $\cl\cherry$ has at most $4$ vertices.
\end{proposition}

\begin{lemma} \label{lem:min-forb-switching}
    Any minimal forbidden subgraph for $\gen\cherry$ has at most $8$ vertices.
\end{lemma}

\begin{proof}
    Let $G$ be a minimal forbidden subgraph for $\gen\cherry$ on $n$ vertices. Assume for the sake of contradiction that $n \ge 9$. Fix a vertex $v$ of $G$. Minimality means $G - v$ is in $\gen\cherry$. By applying a switching to $G$, we may assume that $G - v$ is in $\cl\cherry$, and that $v$ has degree at most $n/2$.
    
    We claim that $G - u$ is in $\cl\cherry$ for every $u \in V(G)$. Take an arbitrary $u \in V(G)$. We may assume that $u \neq v$. Since $G - u \in \gen\cherry$, there exists $U \subseteq V(G)\setminus \sset{u}$ such that $(G-u)^U \in \cl\cherry$. We may assume that $v \not\in U$ by taking the complement of $U$ as a subset of $V(G) \setminus \sset{u}$ in case $v \in U$. Since $G - v \in \cl\cherry$, clearly $G - \sset{u,v}$ is in $\cl\cherry$. Since $(G - \sset{u,v})^{U}$ is also in $\cl\cherry$, by \cref{lem:uniq}, we have that $U$ is trivial as a subset of $V(G) \setminus \sset{u,v}$. Suppose for a moment that $U = V(G) \setminus \sset{u,v}$. Since the degree of $v$ in $(G - u)^{\sset{v}}$ is at least $n / 2 - 1$, the graph $(G - u)^{\sset{v}} \not\in \cl\cherry$, which contradicts the fact that $(G - u)^{\sset{v}} = (G - u)^U \in \cl\cherry$. Therefore $U = \varnothing$, and $G - u = (G - u)^U \in \cl\cherry$.
    
    Based on the claim, we observe that $G$ is a minimal forbidden subgraph for $\cl\cherry$ on at least $9$ vertices, which contradicts \cref{lem:min-forb}.
\end{proof}

\end{document}
